\section{数据连接与公开时点}
\subsection{研究范围及公司日主键}
日收益实值覆盖2015年1月5日至2024年7月31日，共9,091,641行。有效标签覆盖2017年2月6日至2024年7月24日。起始年份的差别来自历史估计和已结束阈值窗口的要求；结束日期的差别来自未来五市场日结果的需要。因子和公告表存在更晚记录，并不能据此延长风险测试。模型范围由收益结果可观测性决定，辅助资料再按公开时点查询。

日收益表中的证券代码和日期组成主键，当前该键已做唯一性核查。ComCd可以参与跨表连接，但不是证券交易代码；CompanyCode又属于另一个命名空间。旧转换器曾随工作簿列顺序选取两类公司码，造成编码混用。修正后，9,091,641行中4,318,444行仅公司码字段变化，收益、日期等标签依赖字段保持不变。这一发现说明来源解析错误和标签算法错误应分别检查，不能一发现公司码问题就宣布全部收益标签失效。

\reportfigure{公司、证券与股东所属证券的不同编码空间}{fig:namespace}{\FigNamespace}{字段连接设计；两命名空间独立保留。观察对应关系仅为候选，不是历史上市有效区间认证。}

图\ref{fig:namespace}中的股东所属证券编码还需要特别处理。基金或证券产品的SecuCd描述的是股东一侧，不是被持股公司的A股代码。若把它用来连接收益表，会产生结构上看似完整、实际指向错误对象的关系。历史证券对应应有有效日期和来源凭证，不能用最新全称或当前股票代码回填全部历史。

\subsection{日期级信息的查询规则}
设市场日历为有序集合$\mathcal T$。对只有公开日期$P_e$的事件，定义生效日
\begin{equation}
 v_e=\min\{t\in\mathcal T:t>P_e\}.
\end{equation}
在预测日$t$读取$v_e\le t$的版本，保存报告期作为业务解释字段，不将其当作生效时点。公告正文、原生数值和后续更正共同进入版本记录；后续更新不能修改首次日已经看到的内容。诉讼首次公开日适合描述事件开始，当前进展则按更新公开日进入控制变量。

公司最近事件年龄取市场日序号差，而不是自然日差。这样周末和长假不会在模型中变成额外收益步长。市场日历来自全市场日期，不来自该公司自己的非缺失收益日；停牌公司也沿同一日历查询事件。每个来源返回原始公开日、生效日、版本键和可用状态，装配器逐一验证最大可用日不晚于预测日。

\section{短期条件尾部标签}
\subsection{起点估计与未来结果}
令$r_{i,s}$为公司$i$日收益，$r_{f,s}$为无风险收益，$F_s$为流通市值加权市场溢酬、SMB和HML。先筛选交换/市场组合$2|A$，避免同一天不同因子组合重复进入连接。在$t$以前最近252个有限有效观测组成集合$\mathcal H_{i,t}$，估计
\begin{equation}
 r_{i,s}-r_{f,s}=\alpha_{i,t}+\beta_{i,t}^{\mathsf T}F_s+\epsilon_{i,s},\quad s\in\mathcal H_{i,t}.
\end{equation}
截距与三因子系数在定义该起点结果时冻结。未来结果为
\begin{equation}
 C_{i,t}^{(5)}=\sum_{k=1}^{5}\left[r_{i,t+k}-r_{f,t+k}-\widehat\alpha_{i,t}-\widehat\beta_{i,t}^{\mathsf T}F_{t+k}\right],
\end{equation}
其中$t+k$表示市场日历中的第$k$个后续交易日，不是第$k$个个股有效记录。未来五日任一必需观测缺失时，不产生该起点标签。

三因子调整使任务侧重公司相对共同市场风险的异常变化，但并不把残差解释为已识别的纯公司冲击。回归窗口、因子估计与遗漏变量都会影响结果。因此，实验使用同一标签比较模型，而不把年度崩盘研究或普通涨跌研究的指标直接用作精度参照。

\subsection{仅使用已结束历史窗口的阈值}
对每个历史起点$u$，其五日结果同样使用$u$时点冻结的系数。预测$t$的阈值只从结果结束日不晚于$t$的窗口选取最近252个有效值，记为$\mathcal C_{i,t}$。排序后固定取
\begin{equation}
 q_{.10,i,t}=C_{i,t,(26)},\quad q_{.05,i,t}=C_{i,t,(13)},\qquad
 Y_{i,t}^{(.10)}=\mathbf1\{C_{i,t}^{(5)}<q_{.10,i,t}\}.
\end{equation}
严格小于阈值才为正类。历史窗口不应以未来重新估计的系数回算；全样本分位数也不能代替起点历史分位数。图\ref{fig:label}呈现了系数和窗口的两条时间路径。

\reportfigure{起点系数、已结束窗口与五日结果的标签构造}{fig:label}{\FigLabel}{标签算法设计及已实施的计算口径；图中未来路径只属于离线监督，不进入预测输入。}

有效历史回归观测和有效结果窗口是两种要求。公司可能已经有252个可回归日，却还没有252个已结束五日窗口；某些窗口还会因停牌或缺失失效。阈值和结果都满足条件才形成共享目标样本。实际P10正类率不需要恰好等于10\%，因为未来分布会变化，窗口之间也有重叠。本课题在每折报告真实正类率，不用名义分位数替代样本统计。

\section{60市场日序列与训练期标准化}
序列主字段为日收益、超额收益和三个因子，得到$X_{i,t}\in\mathbb R^{60\times5}$。相同形状的$M_{i,t}$表示各字段是否有效。缺失项先填0，再与掩码拼接为十通道输入：
\begin{equation}
 \widetilde X_{i,t}=\big[X_{i,t}\odot M_{i,t};M_{i,t}\big],\quad
 h_{i,t}=\operatorname{GRU}_{64}(\widetilde X_{i,t}),\quad
 s_{i,t}=\tanh(W_sh_{i,t}+b_s).
\end{equation}
掩码让模型区分真实零收益与缺失占位。若向前找60个个股有效观测，长期停牌公司的窗口可能跨越远多于60个市场日，序列尺度将不再一致。表\ref{tab:windowcases}给出处理差异。

\reporttable{收益窗口中不同情况的处理}{tab:windowcases}{
\begin{tabularx}{\linewidth}{p{3cm}YY}\toprule
情况&历史输入&标签结果\\\midrule
有效零收益&值0，掩码1&按未来结果完整性计算\\
预测日收益缺失&缺失占位0，掩码0&其他条件满足时可保留目标\\
历史窗口部分缺失&保留相同市场日位置及字段掩码&回归历史按有效观测另行筛选\\
未来五日缺必需观测&未来不参与输入&不产生该起点标签\\
非市场交易日&不占收益序列位置&事件延到下一个市场日生效\\\bottomrule
\end{tabularx}}

各折标准化使用训练期有效值估计均值和尺度，验证与测试只应用已有参数。缺失填充值、异常截断界及正文PCA同样纳入版本记录。不能为方便把2017—2024全部数据一起拟合预处理。基础收益统计对照也遵守这一约定，避免深度模型和强树在不同信息条件下竞争。

GRU提供有序历史表示\cite{r13}，并不自动优于滚动统计。仅收益模型只使用其主BCE，不接收公告、图辅助监督或诉讼控制；同信息融合模型的主干则接收共同的诉讼结构变量。这一区分需要写进输入合同，不能只根据模型名称判断它实际看到了什么。

\section{预告金标：方向、原因和证据如何分开}
\subsection{语义目标与可执行规范}
标注对象是预告正文所表达的经营/财务变化。方向取改善、恶化、基本稳定、混合、不明确；原因是原文明确给出的自由短语数组；证据是1—3段逐字片段。没有原因时记空数组，不反过来把方向改成不明确。未填写、语义不明确和未提及原因分别属于进度状态、语义类别和原因状态。

表\ref{tab:annotationexamples}使用虚构短句说明这一区别。例句不代表真实公司，也不用于模型训练。对标注者而言，主要困难不是编写JSON，而是区分“仍亏损”与“变化改善”，以及判断一段话究竟陈述原因还是重复结果。原因字段不预设行业、市场或管理层分类，以免逼迫标注者补充正文没有给出的解释。

\reporttable{三字段标注的说明用例（全部虚构）}{tab:annotationexamples}{
\begin{tabularx}{\linewidth}{Yp{1.7cm}Y}\toprule
说明文本&方向&原因与证据处理\\\midrule
预计净利润同比增长，主要由于产品销量增加。&改善&原因：产品销量增加；证据复制原句\\
仍为亏损，但亏损额较上年减少。&改善&未提明确原因，原因数组为空；保留原句\\
预计盈利，但净利润比上年下降。&恶化&绝对盈利不能代替同比方向；原因按原文有无填写\\
净利润增长，但主营业务利润下降。&混合&重要指标明确反向，证据同时覆盖两者\\\bottomrule
\end{tabularx}}

\subsection{锁定样本、独立双标与裁决}
主样本1500条按2017—2019训练800、2020验证200、2021—2024年7月测试500固定。200试标是训练子集；300双标是主样本子集，二者交集42。先由A标200、B独立标42，在原稿上计算一致率及方向Kappa，再讨论和冻结规范。余下A1300与B258按同一版本继续，避免试标反馈来自验证/测试而改变规则。

\reportfigure{固定成员下的独立试标、规范冻结与金标放行}{fig:annotation}{\FigAnnotation}{已建立的标注工作流程。当前工作簿为空，流程不等于1500金标已完成；数量均为同一锁定主样本及其子集。}

家族ID、正文和公开日期随来源指纹保留，导入不能增加、删除或修改样本来凑配额。裁决记录绑定独立A/B原稿哈希，讨论后的统一答案不能重算成“原始一致性”。格式预检检查方向枚举、证据数量和正文逐字支持，人工审阅仍负责语义是否正确。只有完成裁决和明确放行的整批样本，才可以导出训练数据。

\section{QLoRA抽取器和正文表示}
\subsection{参数适配及独立质量评价}
Qwen3-8B采用NF4四比特量化和低秩适配，rank16、alpha32、dropout0.05，batch1、梯度累计16。LoRA和QLoRA分别提供低秩适配与量化微调方法\cite{r21,r22}，Qwen3提供中文语义基座\cite{r23}。完整序列上限2048 token，包含系统、用户和答案；正文预算从1024 token起测，不能把全部上限理解为可用正文长度。

截断或跳过超长记录会改变被评价语料，必须记录数量、类别、公开年份和失败原因。测试集只评价冻结模型，不能根据测试表现反复修改提示词或标注规则。抽取评价分别比较原生字段/规则、零样本和QLoRA，报告方向Macro-F1、原因短语评价、证据支持、严格JSON率及清洗后可解析率。解析合法率是接口质量，不是语义正确率。

原因和证据都是自由字符串，不能仅以字符串完全相等判定语义质量。方向可按类别统计；原因需检查是否概括了明确因由、是否混入结果；证据需检查覆盖和语义支持。模型自报置信度不作为正确率。独立测试的误判案例保留在私有记录中，用于解释下游分组结果，不把测试金标直接送入风险模型。

\subsection{正文池化、结构字段与缓存}
对每条去重公告只编码一次，正文向量为
\begin{equation}
 v_e=\frac{\sum_{q}m_q^{body}H_q^{last}}{\max(1,\sum_qm_q^{body})}.
\end{equation}
正文位置掩码排除系统提示、公司包装、padding和生成答案。若把生成JSON也池化，风险表示便混入模型自己生成的解释，无法清楚区分正文信息与生成反馈。图\ref{fig:pooling}说明池化范围。

\reportfigure{正文token与生成答案的池化范围分离}{fig:pooling}{\FigPooling}{拟议的冻结表示生成规则。正文预算和位置按分词器、模板版本验证，不展示真实公告内容。}

事件缓存同时保留预测三字段及原生数值。风险向量的数值结构包括方向编码、原因是否明确、证据可追溯状态和原生字段/缺失标记；自由原因短语与证据原文保留给解释和独立抽取评价，不把它们硬映射为预设行业类别。正文向量承担连续语义信息，风险投影从训练数据学习，不冻结随机矩阵。

\reportfigure{事件缓存键与日查询的版本关系}{fig:cache}{\FigCache}{缓存设计。每一输入或模型组成改变都会形成不同键，向量与预测字段使用同一可追溯版本。}

公司日事件聚合最多使用最近60市场日内的8条有效预告。对投影向量$\widetilde e_j$，内容分数经过数值截断后乘可用性、可追溯质量和年龄衰减，再归一化：
\begin{equation}
 b_j=\exp\{\operatorname{clip}(u^T\widetilde e_j,-20,20)\}m_jq_j e^{-r_eA_j},\quad
 e_i=\frac{\sum_{j\in\mathcal E_i(t)}b_j\widetilde e_j}{\max(\varepsilon,\sum_jb_j)}.
\end{equation}
全无有效事件时返回零表示并令$M_e=0$。质量$q_j$表达格式和来源可追溯状态，不等于人工语义正确概率；公司级$Q_e$按参与聚合的有效事件质量均值形成。若只使用二值质量，它主要起过滤作用，不能把$Q_e=1$解释为语义绝对正确。

诉讼5/20/60日计数、类型和最近年龄通过独立结构路径进入同信息模型。业绩预告缺失不应屏蔽诉讼事实；更新进展也不能回填首次公开日。由于Qwen预训练晚于回测期，本课题评价的是现代抽取器处理历史资料的回顾性效用。原生字段、仅预测结构和匿名文本控制有助于分析记忆风险，但无法证明预训练不存在污染。
